\documentclass[11pt]{article}
\usepackage[margin=1in]{geometry}
\usepackage{enumitem}
% =============================================================================
% Preambles/header.tex — shared LaTeX/Beamer preamble for this project
%
% Use from a Beamer lecture by prepending:
%     \input{header}
% and compiling with TEXINPUTS=../Preambles:$TEXINPUTS (the /compile-latex
% skill does this automatically).
%
% PALETTE CONTRACT. Color names defined here MUST match the names in
% Quarto/theme-template.scss so Beamer and Quarto renderings use the same
% palette. The scripts/check-palette-sync.sh script enforces this.
%
% To customize: edit the HEX values below AND the matching $variable in
% Quarto/theme-template.scss, then run scripts/check-palette-sync.sh.
% =============================================================================

% -----------------------------------------------------------------------------
% Engine detection + encoding
%
% Our /compile-latex skill uses XeLaTeX, where inputenc is unnecessary and
% can produce encoding warnings. Only load inputenc under pdfTeX.
% -----------------------------------------------------------------------------
\usepackage{iftex}
\ifPDFTeX
  \usepackage[utf8]{inputenc}
\fi

\usepackage{amsmath, amssymb, amsthm}
\usepackage{booktabs}
\usepackage{graphicx}

% -----------------------------------------------------------------------------
% PALETTE — mirrors Quarto/theme-template.scss variable names.
% Load xcolor and define colors BEFORE anything that references them
% (notably hyperref, hypersetup, and the Beamer color assignments).
% -----------------------------------------------------------------------------
\usepackage{xcolor}

\definecolor{primary-blue}{HTML}{012169}      % headings, accents
\definecolor{primary-gold}{HTML}{B9975B}      % emphasis, borders
\definecolor{highlight-yellow}{HTML}{F2A900}  % markers, alerts
\definecolor{light-bg}{HTML}{E8EDF5}          % title / section backgrounds
\definecolor{jet}{HTML}{1A1A1A}               % body text

% Semantic colors (used in diagrams and annotations)
\definecolor{positive}{HTML}{15803D}          % good / observed / identified
\definecolor{negative}{HTML}{B91C1C}          % bad / problematic / bias
\definecolor{neutral}{HTML}{525252}           % reference / context

% Highlight accents (match .hi-* classes in the SCSS)
\definecolor{hi-slate}{HTML}{314F4F}
\definecolor{hi-green}{HTML}{2E8B57}
\definecolor{hi-red}{HTML}{C41E3A}

% Semantic aliases so skill templates and snippets can write
% \textcolor{accent}{...} without hard-coding "primary-blue".
\colorlet{accent}{primary-blue}
\colorlet{accent2}{primary-gold}

% -----------------------------------------------------------------------------
% Hyperref — loaded AFTER xcolor + palette so link colors resolve.
% -----------------------------------------------------------------------------
\usepackage{hyperref}
\hypersetup{colorlinks=true, linkcolor=primary-blue, urlcolor=primary-blue,
            citecolor=primary-blue, pdfborder={0 0 0}}

% -----------------------------------------------------------------------------
% Course code — set once per deck (after \input{header}, before \begin{document})
% via \coursecode{CS401}. Shown in the footer on every slide so a deck's course
% is identifiable without opening the title page. Empty by default.
% -----------------------------------------------------------------------------
\makeatletter
\newcommand{\coursecode}[1]{\gdef\@coursecodetext{#1}}
\gdef\@coursecodetext{}
\makeatother

% -----------------------------------------------------------------------------
% Beamer theme assignments (only apply under Beamer).
%
% We detect Beamer via \@ifclassloaded{beamer}, which is the documented,
% non-internal way. This must be wrapped in \makeatletter/\makeatother
% because \@ifclassloaded contains an @-catcode character.
% -----------------------------------------------------------------------------
\makeatletter
\@ifclassloaded{beamer}{%
  \usefonttheme{professionalfonts}
  \setbeamercolor{structure}{fg=primary-blue}
  \setbeamercolor{frametitle}{fg=primary-blue}
  \setbeamercolor{title}{fg=primary-blue}
  \setbeamercolor{subtitle}{fg=primary-gold}
  \setbeamercolor{author}{fg=primary-blue}
  \setbeamercolor{institute}{fg=neutral}
  \setbeamercolor{date}{fg=primary-gold}
  \setbeamercolor{itemize item}{fg=primary-blue}
  \setbeamercolor{itemize subitem}{fg=primary-gold}
  \setbeamercolor{alerted text}{fg=primary-gold}
  \setbeamercolor{block title}{fg=white, bg=primary-blue}
  \setbeamercolor{block body}{bg=light-bg}
  % Semantic block variants: block=definition/concept (blue), exampleblock=
  % worked example (green/positive), alertblock=key takeaway or caution (gold).
  % Keep to <=2 colored boxes per slide across ALL three variants (INV-7).
  \setbeamercolor{block title example}{fg=white, bg=positive}
  \setbeamercolor{block body example}{bg=light-bg}
  \setbeamercolor{block title alerted}{fg=white, bg=primary-gold}
  \setbeamercolor{block body alerted}{bg=light-bg}

  % Minimal footer: course code (left) + page X / N (right), no navigation clutter
  \setbeamertemplate{navigation symbols}{}
  \setbeamertemplate{footline}{%
    \color{neutral}\footnotesize\hspace{0.7em}\@coursecodetext\hfill
    \insertframenumber/\inserttotalframenumber\hspace{0.7em}\vspace{0.3em}}
}{}
\makeatother

% -----------------------------------------------------------------------------
% TikZ setup — libraries the snippet gallery uses
% -----------------------------------------------------------------------------
\usepackage{tikz}
\usetikzlibrary{arrows.meta, positioning, calc, decorations.pathreplacing,
                fit, shapes.geometric, backgrounds}

% Shared TikZ styles — snippets and hand-written diagrams can pick these up
% via `\begin{tikzpicture}[every node/.style={...}]` or by naming specific
% styles (e.g., `\node[dag-node] (X) at (0,0) {X};`).
\tikzset{
  % Node styles for causal DAGs and similar
  dag-node/.style = {
    draw=primary-blue, fill=light-bg, rounded corners=2pt,
    minimum width=1.2cm, minimum height=0.9cm, align=center, font=\small
  },
  decision-node/.style = {
    draw=primary-gold, fill=white, diamond, aspect=1.8,
    minimum width=1.8cm, minimum height=1.2cm, align=center, font=\small,
    inner sep=1pt
  },
  % Edge styles — observed vs counterfactual
  observed-edge/.style = {-{Latex[length=2.5mm]}, thick, draw=jet},
  counterfactual-edge/.style = {-{Latex[length=2.5mm]}, thick, draw=neutral, dashed},
  confound-edge/.style = {-{Latex[length=2.5mm]}, thick, draw=negative, dashed},
  % Data-point styles for plots
  observed-dot/.style = {fill=primary-blue, draw=primary-blue, circle, inner sep=1.2pt},
  counterfactual-dot/.style = {fill=white, draw=neutral, circle, inner sep=1.2pt}
}

% -----------------------------------------------------------------------------
% Convenience macros
% -----------------------------------------------------------------------------
\newcommand{\muted}[1]{\textcolor{neutral}{#1}}
\newcommand{\key}[1]{\textcolor{primary-gold}{\textbf{#1}}}
\newcommand{\good}[1]{\textcolor{positive}{#1}}
\newcommand{\bad}[1]{\textcolor{negative}{#1}}

% Standout frame macro: full-bleed dark background for section transitions.
% Beamer-only (uses \begin{frame}/\setbeamercolor/\usebeamerfont) -- do not
% call from an article-class file (e.g. Notes/<CODE>/*-notes.tex). \muted,
% \key, \good, \bad, and \coursecode above are plain xcolor/gdef and are
% safe to use from either class; verified when Notes/article-class support
% was added.
% Expand: \transitionslide{Your transition title}
\newcommand{\transitionslide}[1]{%
  \begingroup
  \setbeamercolor{background canvas}{bg=primary-blue}
  \begin{frame}[plain]
    \centering
    \vfill
    {\usebeamerfont{title}\color{white}\Large #1\par}
    \vfill
  \end{frame}
  \endgroup
}

% Section divider frame (Beamer-only), an alternative to \transitionslide with a
% darker two-line style: near-black (jet) background, a small white label line
% above a larger gold title, and a thin gold rule along the bottom edge.
% Expand: \sectiondivider{Section 2}{Pointers}
\newcommand{\sectiondivider}[2]{%
  \begingroup
  \setbeamercolor{background canvas}{bg=jet}
  \begin{frame}[plain]
    \vspace*{3.0cm}
    {\color{white}\ttfamily\large #1\par}
    \vspace{0.5cm}
    {\color{primary-gold}\LARGE #2\par}
    \begin{tikzpicture}[remember picture, overlay]
      \draw[primary-gold, line width=2.5pt]
        ($(current page.south west)+(0,0.1)$) -- ($(current page.south east)+(0,0.1)$);
    \end{tikzpicture}
  \end{frame}
  \endgroup
}

% -----------------------------------------------------------------------------
% Channel watermark — "Concepts That Click" (YouTube).
%
% Article-class files (CompetitiveExam Q/A sets, Notes, Assignments) get a
% light diagonal draft watermark on every page plus a centered footer naming
% the channel. Beamer decks get a subtle TikZ background watermark instead
% (the footline already carries the course code). Centralized here so every
% MCQ PDF is branded without per-file edits.
% -----------------------------------------------------------------------------
\makeatletter
\@ifclassloaded{article}{%
  \usepackage{draftwatermark}
  \SetWatermarkText{Concepts That Click}
  \SetWatermarkScale{1}
  \SetWatermarkFontSize{2cm}
  \SetWatermarkLightness{0.86}
  \SetWatermarkAngle{45}
  \usepackage{fancyhdr}
  \pagestyle{fancy}
  \fancyhf{}
  \fancyfoot[C]{\footnotesize\textcolor{neutral}{Concepts That Click~|~YouTube: \href{https://www.youtube.com/@concepts.thatclick}{@concepts.thatclick}}}
  \renewcommand{\headrulewidth}{0pt}
  \renewcommand{\footrulewidth}{0pt}
  \fancypagestyle{plain}{%
    \fancyhf{}%
    \fancyfoot[C]{\footnotesize\textcolor{neutral}{Concepts That Click~|~YouTube: \href{https://www.youtube.com/@concepts.thatclick}{@concepts.thatclick}}}%
    \renewcommand{\headrulewidth}{0pt}%
    \renewcommand{\footrulewidth}{0pt}%
  }
}{}
\@ifclassloaded{beamer}{%
  \setbeamertemplate{background}{%
    \begin{tikzpicture}[remember picture, overlay]
      \node[rotate=45, opacity=0.055, align=center]
        at (current page.center)
        {\Huge\textcolor{neutral}{Concepts That Click}};
    \end{tikzpicture}%
  }
}{}
\makeatother 
\coursecode{JKSSB}
\renewcommand{\thesection}{2.\arabic{section}}
\newtheorem{example}{Example}[section]
\newenvironment{definitionbox}[1]{%
  \par\noindent\textbf{Definition (#1).}\ \itshape
}{\par\vspace{0.3em}}
\title{Types of Computers --- Detailed Lecture Notes}
\author{JKSSB Computer Awareness}
\date{\today}

\begin{document}
\maketitle
\noindent\textbf{Video lesson:}
\href{https://youtu.be/d6zL3YS69kM}{Computer Classifications for JKSSB}

\section{Choose the classification axis first}
Computer labels are useful only when the basis of classification is clear.
``Digital'' describes how data are represented; ``microcomputer'' describes
a traditional size/use class; ``laptop'' describes a form factor. These
labels do not compete. One laptop can correctly be described using several
axes at once.

\begin{center}
\begin{tabular}{p{0.22\textwidth}p{0.34\textwidth}p{0.34\textwidth}}
\toprule
\textbf{Axis} & \textbf{Question} & \textbf{Clue / example}\\
\midrule
Signal representation & How are values represented? & Continuous signal: analog; distinct states: digital\\
Scale and workload & What kind of work or volume is handled? & Many transactions: mainframe; very large calculations: supercomputer\\
Purpose & How broad is the intended task range? & Many different programs: general-purpose; narrow task: special-purpose\\
Integration & Where is the computer installed? & Built into a larger product: embedded\\
Form factor & What physical form does it take? & Desktop, laptop, tablet, or smartphone\\
\bottomrule
\end{tabular}
\end{center}

\begin{definitionbox}{Classification axis}
The feature used to group computers, such as signal representation,
capability/workload, purpose, or physical form.
\end{definitionbox}

Before choosing an answer, identify the question being asked:
\begin{enumerate}
  \item Is the stem about continuous or discrete signals?
  \item Is it about personal use, transaction volume, or very large
    calculations?
  \item Is it about a computer's general purpose or its built-in role in a
    larger product?
  \item Is it naming a shape, such as desktop, laptop, or tablet?
\end{enumerate}
Selecting a familiar label without checking the axis is a common source of
wrong answers.

\section{Classification by signal representation}
\begin{definitionbox}{Analog computer}
A computer that represents or processes quantities as continuously varying
signals over a range.
\end{definitionbox}

\begin{definitionbox}{Digital computer}
A computer that represents or processes data as discrete values, commonly
using binary digits 0 and 1.
\end{definitionbox}

\begin{definitionbox}{Hybrid computer}
A system combining analog measurement or input with digital processing,
control, or display.
\end{definitionbox}

An analog quantity can take intermediate values within a range. For
example, a moving pointer can indicate many positions along a speed range,
not just a fixed list of allowed readings. A digital system represents
information using distinct states; modern computers commonly use binary
states, written as 0 and 1. The important contrast is continuous versus
discrete representation, not old versus new or mechanical versus electronic.

\begin{center}
\begin{tabular}{p{0.18\textwidth}p{0.35\textwidth}p{0.35\textwidth}}
\toprule
\textbf{Class} & \textbf{Representation} & \textbf{Recognition clue}\\
\midrule
Analog & Continuously varying quantity & A value can lie anywhere within a range\\
Digital & Separate symbols or states & A value is encoded using distinct states, commonly 0 and 1\\
Hybrid & Uses both forms at different stages & Measurement or input is analog; processing or display is digital\\
\bottomrule
\end{tabular}
\end{center}

A hybrid arrangement can sense a physical signal and then use digital
circuitry to process and present the measurement. Its classification
describes the complete path, not just the display or one internal component.

\begin{example}
A digital patient monitor may measure a changing physical signal through a
sensor. The sensor/front end handles an analog quantity; conversion allows
digital processing; the screen shows a numeric result. When a question asks
about the complete measurement-and-processing arrangement, ``hybrid'' is
the useful classification. The monitor's display alone does not determine
the signal class of its entire system.
\end{example}

\section{Traditional size and capability classes}
\begin{definitionbox}{Microcomputer}
A personal-use computer built around a microprocessor; common examples
include desktop computers, laptops, tablets, and smartphones.
\end{definitionbox}

\begin{definitionbox}{Minicomputer}
A historical term for a mid-range computer class, often described as
supporting multiple users; it is less common in current product marketing.
\end{definitionbox}

\begin{definitionbox}{Mainframe}
A computer class associated with reliable, high-volume input/output and
large numbers of concurrent users or transactions.
\end{definitionbox}

\begin{definitionbox}{Supercomputer}
A high-performance computer designed for very large computational
workloads, often using parallel processing.
\end{definitionbox}

\begin{center}
\begin{tabular}{p{0.2\textwidth}p{0.34\textwidth}p{0.34\textwidth}}
\toprule
\textbf{Class} & \textbf{Primary clue} & \textbf{Example / caution}\\
\midrule
Microcomputer & Personal-use computer built around a microprocessor & Desktop, laptop, tablet, or smartphone; not defined only by its outer size\\
Minicomputer & Historical mid-range, multi-user class & A traditional category; less common in current product marketing\\
Mainframe & High transaction and input/output throughput for many users & Large banking or account-processing workload\\
Supercomputer & Very high computational performance, often parallel & Large climate or engineering simulation\\
\bottomrule
\end{tabular}
\end{center}

Mainframes and supercomputers may both be powerful. The distinction in an
exam question is usually the workload. A mainframe is associated with
throughput, reliability, and many transactions or users. A supercomputer is
associated with computationally intensive scientific or engineering
problems. ``Throughput'' here means the amount of work completed over time;
the mainframe clue is handling many transactions reliably, rather than
solving one unusually large numerical problem. Parallel processing divides
computational work so multiple processing units can perform parts of it at
the same time. Physical size alone is not a dependable classification rule.

\begin{example}
A bank processes a large number of account transactions for many customers.
The emphasis is concurrent transaction throughput and service reliability,
so the mainframe class is the best match. A research centre running a
large-scale climate simulation emphasizes high-performance computation, so
the supercomputer class is the better match.
\end{example}

\begin{example}
A question describes a ``mid-range, multi-user computer'' and asks for its
traditional category. The intended answer is minicomputer. Do not infer that
the label is common in modern consumer advertisements; it is a historical
classification term.
\end{example}

\section{Classification by purpose and form}
\begin{definitionbox}{General-purpose computer}
A computer capable of performing many different tasks by running different
programs.
\end{definitionbox}

\begin{definitionbox}{Special-purpose computer}
A computer designed primarily for a narrower task or a limited set of
related tasks.
\end{definitionbox}

\begin{definitionbox}{Embedded computer}
A computer built into a larger device to control or support that device's
function.
\end{definitionbox}

\begin{center}
\begin{tabular}{p{0.2\textwidth}p{0.37\textwidth}p{0.3\textwidth}}
\toprule
\textbf{Term} & \textbf{What it describes} & \textbf{Example}\\
\midrule
General-purpose & Broad range of tasks through different programs & Personal computer used for documents, browsing, and media\\
Special-purpose & Narrow task or small set of related tasks & Controller dedicated to an appliance function\\
Embedded & Integration into a larger product & Washing-machine controller\\
Form factor & Physical arrangement or shape & Desktop, laptop, tablet, smartphone\\
\bottomrule
\end{tabular}
\end{center}

An embedded computer is commonly special-purpose, but these terms name
different aspects: ``embedded'' describes where the computer is integrated;
``special-purpose'' describes the range of tasks it is intended to perform.
A desktop computer that runs office, browsing, and media programs is a
general-purpose system. A washing-machine controller is built into a larger
appliance and handles a limited control task. A computer can therefore be
both embedded and special-purpose; the terms are compatible, not competing
answers. In contrast, ``laptop'' says how a computer is physically arranged
and does not by itself tell whether the data representation is analog or
digital.

\begin{example}
A smartphone can be classified as digital by signal representation, a
microcomputer by personal-use class, and a mobile device by form factor.
These labels can all be true at the same time because they answer different
questions.
\end{example}

\begin{example}
A washing machine's controller is embedded in the appliance and is
special-purpose because it manages a limited set of appliance functions.
Calling it a mainframe because it controls a whole machine confuses physical
integration with workload class.
\end{example}

\section{A reliable matching procedure}
When an exam stem describes a device, underline the clue that indicates the
axis. ``Continuous'' points to analog; ``discrete'' or binary points to
digital; ``sensor plus digital processing'' may point to hybrid. ``Many
concurrent transactions'' suggests mainframe. ``Large scientific
calculations'' suggests supercomputer. ``Built into an appliance'' points to
embedded. ``Desktop'' or ``laptop'' identifies form factor, not signal type.

\begin{example}
Classify a laptop used by an accounts clerk. By representation it is digital;
by traditional size/use class it is a microcomputer; by form it is a laptop;
and by purpose it is general-purpose if it runs many kinds of programs.
None of these labels invalidates the others.
\end{example}

\subsection*{Computer type and generation}
A computer's type may describe its signal class, scale, purpose, or form.
Generation instead names a historical technology grouping. Keep these axes
separate: integrated circuits are associated with the third generation,
while microprocessors are associated with the fourth.

This distinction prevents a common category error. A laptop can be called a
digital microcomputer and can belong to a particular historical generation
based on the technology used in its development. ``Laptop'' identifies its
form, ``microcomputer'' its traditional user/capability class, ``digital''
its representation, and ``generation'' its place in a historical grouping.
Each answer must match the axis named in the question.

\begin{example}
For a washing-machine controller, ``embedded'' answers where the computer is
installed. ``Special-purpose'' answers the range of work it is designed to
perform. The controller can correctly have both labels at once. If a
question asks how it represents data, that is a separate signal
classification and must be answered from that clue.
\end{example}

\subsection*{Exercises}
\begin{enumerate}[label=\arabic*.]
  \item Classify a smartphone on the signal, traditional class, and form
    axes.
  \item A system is described as handling many simultaneous banking
    transactions reliably. Which class is the strongest match, and why?
  \item Distinguish an embedded computer from a special-purpose computer.
  \item Why is ``laptop'' not an answer to a question asking whether a
    computer is analog or digital?
  \item A digital instrument measures a continuously changing physical
    quantity and processes the result digitally. Explain when ``hybrid'' is
    an appropriate label.
\end{enumerate}

\subsection*{Solutions}
\begingroup\small
\begin{enumerate}[label=\arabic*.]
  \item It is digital by representation, a microcomputer by personal-use
    class, and a smartphone by form factor.
  \item Mainframe. The clues are many simultaneous transactions and
    reliability/throughput, rather than a single extremely large calculation.
  \item Embedded describes integration into a larger product. Special-purpose
    describes a narrow intended task. An embedded controller is often
    special-purpose, but the terms describe different properties.
  \item Laptop names physical form. Analog/digital names the representation
    of data, so the question asks about a different classification axis.
  \item It is appropriate when the complete arrangement combines analog
    sensing or measurement with digital conversion and processing.
\end{enumerate}
\endgroup
\end{document}
